
\documentclass{article}
\usepackage{colortbl}
\usepackage{makecell}
\usepackage{multirow}
\usepackage{supertabular}

\begin{document}

\newcounter{utterance}

\centering \large Interaction Transcript for game `cladder', experiment `full\_v1.5\_default', episode 7517 with refineonly{-}Qwen3.5{-}27B.
\vspace{24pt}

{ \footnotesize  \setcounter{utterance}{1}
\setlength{\tabcolsep}{0pt}
\begin{supertabular}{c@{$\;$}|p{.15\linewidth}@{}p{.15\linewidth}p{.15\linewidth}p{.15\linewidth}p{.15\linewidth}p{.15\linewidth}}
    \# & \multicolumn{2}{c}{Player} && \multicolumn{2}{c}{Game Master} \\
    \hline

    \theutterance \stepcounter{utterance}  
    & & & \multicolumn{4}{p{0.6\linewidth}}{
        \cellcolor[rgb]{0.9,0.9,0.9}{
            \makecell[{{p{\linewidth}}}]{
                \texttt{\tiny{[P1$\langle$GM]}}
                \texttt{You are an expert in causal inference. The following question is not a typical commonsense query, but rather a meticulously designed question created by a professor specializing in causal inference, intended to assess the students' mastery of the course content.} \\
\texttt{Imagine a self{-}contained, hypothetical world with only the following conditions, and without any unmentioned factors or causal relationships: Wibl has a direct effect on vubr and zuph. Zuph has a direct effect on uvzi. Vubr has a direct effect on uvzi. Method 1: We look at how zuph correlates with uvzi case by case according to vubr. Method 2: We look directly at how zuph correlates with uvzi in general. To understand how zuph affects uvzi, is it more correct to use the Method 1 than Method 2?} \\
\texttt{Start your answer with 'Yes' or 'No', followed by additional reasoning or evidence to support your explanation.} \\
            }
        }
    }
     \\ \\

    \theutterance \stepcounter{utterance}  
    & & \multicolumn{4}{p{0.6\linewidth}}{
        \cellcolor[rgb]{0.95,0.95,0.95}{
            \makecell[{{p{\linewidth}}}]{
                \texttt{\tiny{[P1$\leftrightarrow$P]}}
                \texttt{{[}DRAFT{]}} \\
\texttt{Yes} \\
\\ 
\texttt{To determine the correct approach, we must first map the causal structure described in the hypothetical world. The conditions state that **Wibl** causes both **Vubr** and **Zuph**, making Wibl a common cause (confounder) of the two. Furthermore, **Zuph** has a direct effect on **Uvzi**, and **Vubr** also has a direct effect on **Uvzi**.} \\
\\ 
\texttt{This structure creates a specific causal pathway: \$Wibl \textbackslash{}rightarrow Zuph \textbackslash{}rightarrow Uvzi\$ and \$Wibl \textbackslash{}rightarrow Vubr \textbackslash{}rightarrow Uvzi\$. Crucially, there is also a path \$Zuph \textbackslash{}leftarrow Wibl \textbackslash{}rightarrow Vubr \textbackslash{}rightarrow Uvzi\$. In this configuration, **Vubr** acts as a **mediator** of the effect of Wibl on Uvzi, but more importantly for the relationship between Zuph and Uvzi, Vubr is a **collider** on the path between Zuph and Uvzi if we consider the path \$Zuph \textbackslash{}leftarrow Wibl \textbackslash{}rightarrow Vubr \textbackslash{}rightarrow Uvzi\$? No, let's re{-}evaluate the graph structure carefully.} \\
\\ 
\texttt{The graph is:} \\
\texttt{1. \$Wibl \textbackslash{}rightarrow Zuph\$} \\
\texttt{2. \$Wibl \textbackslash{}rightarrow Vubr\$} \\
\texttt{3. \$Zuph \textbackslash{}rightarrow Uvzi\$} \\
\texttt{4. \$Vubr \textbackslash{}rightarrow Uvzi\$} \\
\\ 
\texttt{We are interested in the causal effect of **Zuph** on **Uvzi**.} \\
\texttt{The direct causal path is \$Zuph \textbackslash{}rightarrow Uvzi\$.} \\
\texttt{Is there a backdoor path from Zuph to Uvzi?} \\
\texttt{A backdoor path would start with an arrow pointing *into* Zuph. The only such arrow is from Wibl (\$Wibl \textbackslash{}rightarrow Zuph\$).} \\
\texttt{From Wibl, we can go to Vubr (\$Wibl \textbackslash{}rightarrow Vubr\$), and then to Uvzi (\$Vubr \textbackslash{}rightarrow Uvzi\$).} \\
\texttt{So, the backdoor path is: \$Zuph \textbackslash{}leftarrow Wibl \textbackslash{}rightarrow Vubr \textbackslash{}rightarrow Uvzi\$.} \\
\\ 
\texttt{To estimate the causal effect of Zuph on Uvzi, we must block this backdoor path.} \\
\texttt{{-} **Method 2** (looking at the general correlation) leaves this backdoor path open. Because Wibl affects both Zuph and Vubr (which affects Uvzi), there is a spurious correlation between Zuph and Uvzi induced by Wibl. Therefore, the marginal correlation (Method 2) will be biased; it conflates the direct effect of Zuph with the confounding effect of Wibl.} \\
\texttt{{-} **Method 1** involves conditioning on **Vubr** (looking case by case according to Vubr).} \\
\texttt{  {-} Does conditioning on Vubr} \\
            }
        }
    }
    & \\ \\

    \theutterance \stepcounter{utterance}  
    & \multicolumn{4}{p{0.6\linewidth}}{
        \cellcolor[rgb]{0.561,0.737,0.561}{
            \makecell[{{p{\linewidth}}}]{
                \texttt{\tiny{[P1$\rangle$GM]}}
                \texttt{No} \\
\\ 
\texttt{To determine the correct approach, we must first map the causal structure described in the hypothetical world. The conditions state that **Wibl** causes both **Vubr** and **Zuph**, making Wibl a common cause (confounder) of the two. Furthermore, **Zuph** has a direct effect on **Uvzi**, and **Vubr** also has a direct effect on **Uvzi**.} \\
\\ 
\texttt{The causal graph is:} \\
\texttt{1. \$Wibl \textbackslash{}rightarrow Zuph\$} \\
\texttt{2. \$Wibl \textbackslash{}rightarrow Vubr\$} \\
\texttt{3. \$Zuph \textbackslash{}rightarrow Uvzi\$} \\
\texttt{4. \$Vubr \textbackslash{}rightarrow Uvzi\$} \\
\\ 
\texttt{We are interested in the causal effect of **Zuph** on **Uvzi**.} \\
\texttt{The direct causal path is \$Zuph \textbackslash{}rightarrow Uvzi\$.} \\
\texttt{Is there a backdoor path from Zuph to Uvzi?} \\
\texttt{A backdoor path starts with an arrow pointing *into* Zuph. The only such arrow is from Wibl (\$Wibl \textbackslash{}rightarrow Zuph\$).} \\
\texttt{From Wibl, we can go to Vubr (\$Wibl \textbackslash{}rightarrow Vubr\$), and then to Uvzi (\$Vubr \textbackslash{}rightarrow Uvzi\$).} \\
\texttt{So, the backdoor path is: \$Zuph \textbackslash{}leftarrow Wibl \textbackslash{}rightarrow Vubr \textbackslash{}rightarrow Uvzi\$.} \\
\\ 
\texttt{To estimate the causal effect of Zuph on Uvzi, we must block this backdoor path.} \\
\texttt{{-} **Method 2** (looking at the general correlation) leaves this backdoor path open. Because Wibl affects both Zuph and Vubr (which affects Uvzi), there is a spurious correlation between Zuph and Uvzi induced by Wibl. Therefore, the marginal correlation (Method 2) will be biased; it conflates the direct effect of Zuph with the confounding effect of Wibl.} \\
\texttt{{-} **Method 1** involves conditioning on **Vubr** (looking case by case according to Vubr).} \\
\texttt{  {-} Does conditioning on Vubr block the backdoor path? No. In the path \$Zuph \textbackslash{}leftarrow Wibl \textbackslash{}rightarrow Vubr \textbackslash{}rightarrow Uvzi\$, Vubr is a non{-}collider (specifically, it is a node where arrows meet head{-}to{-}tail: \$Wibl \textbackslash{}rightarrow Vubr \textbackslash{}rightarrow Uvzi\$). Conditioning on a non{-}collider on a path does not block the path; in fact, it is part of the open path.} \\
\texttt{  {-} More critically, conditioning on Vubr does not block the flow of information from Wibl to Zuph. The path remains open: \$Zuph \textbackslash{}leftarrow Wibl \textbackslash{}rightarrow Vubr\$ (conditioned) \$\textbackslash{}rightarrow Uvzi\$. Since Wibl is not conditioned on, the confounding bias persists.} \\
            }
        }
    }
    & & \\ \\

    \theutterance \stepcounter{utterance}  
    & & & \multicolumn{2}{p{0.3\linewidth}}{
        \cellcolor[rgb]{0.9,0.9,0.9}{
            \makecell[{{p{\linewidth}}}]{
                \texttt{\tiny{[GM$|$GM]}}
                \texttt{No} \\
            }
        }
    }
    & & \\ \\

    \theutterance \stepcounter{utterance}  
    & & & \multicolumn{2}{p{0.3\linewidth}}{
        \cellcolor[rgb]{0.9,0.9,0.9}{
            \makecell[{{p{\linewidth}}}]{
                \texttt{\tiny{[GM$|$GM]}}
                \texttt{yes} \\
            }
        }
    }
    & & \\ \\

    \theutterance \stepcounter{utterance}  
    & & & \multicolumn{2}{p{0.3\linewidth}}{
        \cellcolor[rgb]{0.9,0.9,0.9}{
            \makecell[{{p{\linewidth}}}]{
                \texttt{\tiny{[GM$|$GM]}}
                \texttt{game\_result = LOSE} \\
            }
        }
    }
    & & \\ \\

\end{supertabular}
}

\end{document}
